
\documentclass{article}
\usepackage{colortbl}
\usepackage{makecell}
\usepackage{multirow}
\usepackage{supertabular}

\begin{document}

\newcounter{utterance}

\centering \large Interaction Transcript for game `mmlu\_pro', experiment `mmlu\_pro\_psychology', episode 594 with reflexion{-}llama{-}3.1{-}8b{-}instruct.
\vspace{24pt}

{ \footnotesize  \setcounter{utterance}{1}
\setlength{\tabcolsep}{0pt}
\begin{supertabular}{c@{$\;$}|p{.15\linewidth}@{}p{.15\linewidth}p{.15\linewidth}p{.15\linewidth}p{.15\linewidth}p{.15\linewidth}}
    \# & \multicolumn{2}{c}{Player} && \multicolumn{2}{c}{Game Master} \\
    \hline

    \theutterance \stepcounter{utterance}  
    & & & \multicolumn{4}{p{0.6\linewidth}}{
        \cellcolor[rgb]{0.9,0.9,0.9}{
            \makecell[{{p{\linewidth}}}]{
                \texttt{\tiny{[P1$\langle$GM]}}
                \texttt{You have attempted this same task before and did not succeed. These are the plans you wrote after those attempts.} \\
\texttt{Trial \#1: New plan: I got stuck in a loop of answering questions without considering the overall strategy of the game. I was simply answering each question individually without thinking about how the answers would relate to each other or how they would help me achieve the target state of the game. Next time, I will take a step back and think about the overall strategy of the game. I will try to identify patterns and connections between the questions and use that information to inform my answers. I will also try to think about how each answer will contribute to achieving the target state of the game, rather than just answering each question individually.} \\
\\ 
\texttt{Trial \#2: New plan: In my previous attempts, I was focused on answering each question individually without considering the overall strategy of the game. However, I realized that the questions are multiple{-}choice and are likely to be related to each other in some way. Next time, I will take a closer look at the options and try to identify any patterns or connections between the questions. I will also try to use the information from previous questions to inform my answers to subsequent questions. Additionally, I will pay closer attention to the options and try to eliminate any obviously incorrect answers before making a final choice. This should help me to avoid getting stuck in a loop and improve my chances of achieving the target state of the game.} \\
\\ 
\\ 
\texttt{The following are multiple choice questions (with answers) about psychology. Answer only with "The answer is (X)" where X is the correct letter choice.} \\
\\ 
\texttt{Question:} \\
\texttt{Pascale is interested in the processing strategies children use to learn new information. Pascale would best be classified as what type of psychologist?} \\
\texttt{Options:} \\
\texttt{A. social} \\
\texttt{B. school} \\
\texttt{C. sociocultural} \\
\texttt{D. forensic} \\
\texttt{E. behaviorist} \\
\texttt{F. health} \\
\texttt{G. clinical} \\
\texttt{H. cognitive} \\
\texttt{I. psychoanalytic} \\
\texttt{J. developmental} \\
\texttt{The answer is (H)} \\
\\ 
\texttt{Question:} \\
\texttt{According to Caplan's model of consultee{-}centered case consultation, the consultant is primarily interested in} \\
\texttt{Options:} \\
\texttt{A. identifying the causes and solutions of the client's presenting problems} \\
\texttt{B. establishing a hierarchy of authority to enable effective decision making} \\
\texttt{C. ensuring the consultee adheres strictly to a predetermined action plan} \\
\texttt{D. proposing multiple alternative solutions for the consultee to choose from} \\
\texttt{E. identifying the strengths and weaknesses of the consultee's current approach} \\
\texttt{F. presenting a single, well{-}defined and unambiguous course of action for the consultant to overcome skills deficits} \\
\texttt{G. developing a comprehensive treatment plan for the client} \\
\texttt{H. identifying and eliminating the causes of the consultee's difficulties in handling a problem} \\
\texttt{I. focusing on the consultant's personal growth and development} \\
\texttt{J. focusing on the relationship between the client and the consultee} \\
\texttt{The answer is (H)} \\
\\ 
\texttt{Question:} \\
\texttt{According to the Individuals with Disabilities Education Improvement Act, which of the following must an educational agency do before it changes the educational placement of a student with a disability?} \\
\texttt{Options:} \\
\texttt{A. Notify the parents in writing} \\
\texttt{B. Obtain the child's consent} \\
\texttt{C. Obtain a court order} \\
\texttt{D. Conduct a new evaluation of the child's disability} \\
\texttt{E. Discuss with the child's psychologist} \\
\texttt{F. Give the child a trial period in the new environment} \\
\texttt{G. Obtain parental consent} \\
\texttt{H. Notify the local education authority} \\
\texttt{I. Arrange a meeting with all teachers and administrators} \\
\texttt{J. Obtain school board approval} \\
\texttt{The answer is (A)} \\
\\ 
\texttt{Question:} \\
\texttt{Ani believes that her attitudes and behavior play a central role in what happens to her. Such a belief is likely to be associated with} \\
\texttt{Options:} \\
\texttt{A. low self{-}esteem.} \\
\texttt{B. a strong id.} \\
\texttt{C. a high level of anxiety.} \\
\texttt{D. a strong superego.} \\
\texttt{E. high levels of self{-}consciousness.} \\
\texttt{F. an external locus of control.} \\
\texttt{G. an inferiority complex.} \\
\texttt{H. a low level of self{-}awareness.} \\
\texttt{I. low self{-}efficacy.} \\
\texttt{J. an internal locus of control.} \\
\texttt{The answer is (J)} \\
\\ 
\texttt{Question:} \\
\texttt{In terms of Hofstede’s (1980) five cultural dimensions, the United States scores at the top of the scale on:} \\
\texttt{Options:} \\
\texttt{A. individualism and long{-}term orientation.} \\
\texttt{B. individualism and power distance.} \\
\texttt{C. uncertainty avoidance.} \\
\texttt{D. long{-}term orientation.} \\
\texttt{E. individualism.} \\
\texttt{F. individualism and masculinity.} \\
\texttt{G. long{-}term orientation and uncertainty avoidance.} \\
\texttt{H. power distance.} \\
\texttt{I. power distance and masculinity.} \\
\texttt{J. N/A} \\
\texttt{The answer is (E)} \\
\\ 
\texttt{Question:} \\
\texttt{Light enters the eye through an opening in the} \\
\texttt{Options:} \\
\texttt{A. retina.} \\
\texttt{B. cornea.} \\
\texttt{C. iris.} \\
\texttt{D. blind spot.} \\
\texttt{E. lens.} \\
\texttt{F. vitreous humor.} \\
\texttt{G. pupil.} \\
\texttt{H. sclera.} \\
\texttt{I. optic nerve.} \\
\texttt{J. conjunctiva.} \\
            }
        }
    }
     \\ \\

    \theutterance \stepcounter{utterance}  
    & \multicolumn{4}{p{0.6\linewidth}}{
        \cellcolor[rgb]{0.561,0.737,0.561}{
            \makecell[{{p{\linewidth}}}]{
                \texttt{\tiny{[P1$\rangle$GM]}}
                \texttt{Based on the provided information, I will answer the questions with the correct letter choice.} \\
\\ 
\texttt{1. The answer is (H)} \\
\\ 
\texttt{2. The answer is (H)} \\
\\ 
\texttt{3. The answer is (A)} \\
\\ 
\texttt{4. The answer is (J)} \\
\\ 
\texttt{5. The answer is (E)} \\
\\ 
\texttt{6. The answer is (G)} \\
            }
        }
    }
    & & \\ \\

    \theutterance \stepcounter{utterance}  
    & & & \multicolumn{2}{p{0.3\linewidth}}{
        \cellcolor[rgb]{0.9,0.9,0.9}{
            \makecell[{{p{\linewidth}}}]{
                \texttt{\tiny{[GM$|$GM]}}
                \texttt{ParseError: Response contains more than a one answer, but the expected format is to only give one.} \\
            }
        }
    }
    & & \\ \\

    \theutterance \stepcounter{utterance}  
    & & & \multicolumn{2}{p{0.3\linewidth}}{
        \cellcolor[rgb]{0.9,0.9,0.9}{
            \makecell[{{p{\linewidth}}}]{
                \texttt{\tiny{[GM$|$GM]}}
                \texttt{game\_result = ABORT} \\
            }
        }
    }
    & & \\ \\

    \theutterance \stepcounter{utterance}  
    & & \multicolumn{4}{p{0.6\linewidth}}{
        \cellcolor[rgb]{0.95,0.95,0.95}{
            \makecell[{{p{\linewidth}}}]{
                \texttt{\tiny{[P1$\leftrightarrow$P]}}
                \texttt{Post{-}episode Reflexion (not sent to GM):} \\
\texttt{New plan: In my previous attempts, I was focused on answering each question individually without considering the overall strategy of the game. However, I realized that the questions are multiple{-}choice and are likely to be related to each other in some way. Next time, I will take a closer look at the options and try to identify any patterns or connections between the questions. I will also try to use the information from previous questions to inform my answers to subsequent questions. Additionally, I will pay closer attention to the format of the game, which requires me to answer with "The answer is (X)" where X is the correct letter choice. I will make sure to follow this format exactly, without adding any extra text or deviating from it. This should help me to avoid getting stuck in a loop and improve my chances of achieving the target state of the game.} \\
            }
        }
    }
    & \\ \\

\end{supertabular}
}

\end{document}
